% Classical lazy-sampling
\begin{figure}[t]
  \centering
  \begin{pchstack}[center]
    \procedure{$\RO(x)$}{
      \pcif \Db[x] = \bot \pcdo \Db[x] \sample \bin^m \\
      \pcreturn \Db[x]
    }
  \end{pchstack}
  \caption{
    \label{fig:lazy-sampling}
    A lazy-sampled random oracle from $n$ to $m$ bits. 
    Here, $\Db$ is a query database. Note that the oracle is \emph{stateful},
    since the database is maintained between calls.
  }
\end{figure}
%
\section{The Random Oracle Model}\label{sec:rom}
  Random oracles are nice. They are also very easy to simulate, a process known as
    lazy-sampling, as shown in \autoref{fig:lazy-sampling}.

  \subsection{From a One-way PKE to an Indistinguishable KEM: The $\U$-transform}
    The Fujisaki-Okamoto transformation~\cite{C:FujOka99}, which underlies many post-quantum
      KEMs including the now-standardized ML-KEM (Kyber)~\cite{Kyber}, can be
      framed as the composition of two 
      transformations, \T and \U, as defined by Hofheinz
      et al.~\cite{TCC:HofHovKil17} (hereafter HHK). 
      Starting from a public-key
      encryption scheme, the first transformation (\T) derandomizes the
      PKE by having its randomness be derived from a hash of the message (and possibly
      the public key), and additionally adds a re-encryption check to the
      decryption algorithm, while the second transformation builds a KEM from a
      deterministic PKE by sampling a uniform message, encrypting it, and
      hashing it to produce the ephemeral key $\eK$.
  
    In the following, we will focus on this second transformation, $\U$,
      and build our KEM starting with a deterministing PKE that satisfies
      one-way security.\footnote{
        HHK defined several variants of the \U-transform that differ from the variant discussed here
    		in how they handle invalid ciphertexts and in what they hash into the
        output key; our variant corresponds to \UExpMess in their notation.
        %\hhnote{
        %  Should we do the more general \UExpMess transform instead? Or does that
        %  just muddy the message?
        %}
      } (See App.~\ref{sec:ttransform} for details on the \T-transform.)

    \hhnote{
      Leave this version without the extra re-encryption check (which also
      matches the original \U), and define a new \U with the re-encryption check
      when we later want CCA security.
    }
   
    \begin{definition}[\U-transform]
      To a deterministic PKE scheme $\dPKE = (\KG, \Enc, \Dec)$ and a hash
      function $\ROkey: \MSpace \to \KSpace$,
      we associate the KEM $\KEM \coloneqq \U[\PKE, \ROkey]$ defined in \autoref{fig:def:U-transform}.
      \begin{figure}[ht]
   	\begin{center}
   		\fbox{
   			\small
   			\nicoresetlinenr
   			%\begin{minipage}[t]{3.3cm}
   			%	\underline{$\UKG()$}
   			%	\begin{nicodemus}
   			%		\item $(\pk, \sk) \sample \KG()$
   			%		%\item $\rejectionSeed \sample \MSpace$
   			%		%\item $\sk' \leftarrow (\sk, \rejectionSeed)$
   			%		\item \pcreturn $(\pk, \sk)$
   			%	\end{nicodemus}
   			%\end{minipage}
   			%
   			%\quad
   			
   			\begin{minipage}[t]{3cm}
           \underline{$\UEncaps_\pk()$}
   				\begin{nicodemus}
   					\item $m \sample \MSpace$
   					\item $c \gets \Enc(\pk, m)$
   					\item $\eK \gets \ROkey(m)$
   					\item \pcreturn $(\eK, c)$
   				\end{nicodemus}
   			\end{minipage}
   			
   			\quad
   			
   			\begin{minipage}[t]{3.6cm}
          \underline{$\UDecaps_\sk(c)$}
   				\begin{nicodemus}
   					%\item Parse $\sk' := (\sk, \rejectionSeed)$
            \item $m' \gets \Dec_\sk(c)$
   					\item \pcif $m' = \bot \pcthen \pcreturn \bot$
   						%\item \pcreturn $\ROkey(\rejectionSeed, c)$
   					\item \pcreturn $\ROkey(m')$
   				\end{nicodemus}
   			\end{minipage}
   		}
   	\end{center}
   	\caption{
   	  \label{fig:def:U-transform}			
        Algorithms of $\U[\PKE, \ROkey] \eqqcolon (\KG, \UEncaps, \UDecaps)$.
     }
   \end{figure} 
   \end{definition}
  
  \subsection{What Goes Wrong}
    \todoenum[rough structure]{
      \item A discussion about why the previous proofs do not yield
        post-quantum security 
      \item Where the proofs break down, and why they do not admit direct
        lifting (not history-free)
    }
